\documentclass[11pt,a4paper]{article}
\usepackage[T1]{fontenc}
\usepackage[utf8]{inputenc}
\usepackage[polish]{babel}
\usepackage{lmodern}
\usepackage{microtype}
\usepackage[a4paper,margin=2.25cm]{geometry}
\usepackage{amsmath,amssymb,mathtools}
\usepackage{booktabs,longtable,tabularx}
\usepackage{xcolor}
\usepackage{hyperref}
\usepackage{listings}
\usepackage{enumitem}
\hypersetup{colorlinks=true,linkcolor=blue,urlcolor=blue}
\lstset{basicstyle=\ttfamily\small,breaklines=true,frame=single}

\title{Dodatek techniczny\\Generyczny rejestr parametrów jakości powietrza}
\author{Edward Szczypka}
\date{\today}

\begin{document}
\maketitle

\section{Zasada rozdzielenia ról}

Dla każdego parametru $p$ definiuje się niezależne predykaty:
\[
 C_p,\ H_p,\ A_p,\ F_p,\ S_p\in\{0,1\},
\]
gdzie odpowiednio oznaczają pobieranie bieżące, historyczny backfill, rolę
cechy pomocniczej, rolę celu prognozy oraz publikację powierzchni
przestrzennej. W szczególności:
\[
 C_p=1\not\Rightarrow F_p=1.
\]
Pozwala to zbierać i audytować NO2 lub O3 bez automatycznego mnożenia liczby
modeli produkcyjnych.

\section{Kontrakt definicji}

\begin{lstlisting}[language=Python]
class AirParameterDefinition:
    code: str
    aliases: tuple[str, ...]
    canonical_unit: str
    cadence_hours: int
    collect_current: bool
    historical_backfill: bool
    auxiliary_feature: bool
    forecast_target: bool
    spatial_surface: bool
    allow_negative: bool
    valid_min: float | None
    valid_max: float | None
    exceedance_threshold: float | None
    spike_absolute: float | None
    annual_api_indicator: str
    prepared_archive_tokens: tuple[str, ...]
    algorithms: tuple[str, ...]
\end{lstlisting}

Rejestr jest jedynym źródłem prawdy dla kolektorów, importu historycznego,
walidacji, treningu, predykcji i publikacji map.

\section{Aliasowanie}

Funkcja normalizacji usuwa różnice wielkości liter, polskie znaki, spacje,
podkreślenia i zapis przecinkowy. Obsługiwane są również indeksy dolne:
\[
 \mathrm{NO_2}\equiv\texttt{NO2},\qquad
 \mathrm{PM2,5}\equiv\mathrm{PM2.5}.
\]
Nieznany parametr może zostać zachowany tylko w katalogu metadanych albo
pobrany do audytu zgodnie z polityką \texttt{unknown\_sensor\_policy}.

\section{Bridge pobierania}

\begin{lstlisting}
AirParameterRegistry
        |
        +-- collect_current ------> GiosCollector
        +-- historical_backfill --> GiosHistoryImporter
        +-- RangeAwareBackfillBridge
        +-- local / object_store / hybrid cache
\end{lstlisting}

CLI bez jawnej listy wybiera parametry przez rolę. Zapis \texttt{ALL} oznacza
wszystkie parametry danej roli, a nie każdy kod znaleziony przypadkowo w
metadanych.

\section{Generyczne cechy pomocnicze}

Dla celu $p$ i pomocniczego parametru $q\ne p$ w chwili bazowej $t$ tworzy
się:
\[
 z_q(t)=\left(q_t,q_{t-1},q_{t-6},q_{t-24}\right).
\]
Łączenie następuje po stacji i dokładnej godzinie. Nie wykorzystuje się
$q_{t+h}$, dlatego konstrukcja nie powoduje wycieku informacji z przyszłości.

\section{Snapshot i publiczna aplikacja}

Snapshot zawiera bieżące wartości parametrów aktywnych w co najmniej jednej
roli operacyjnej. Manifest przestrzenny publikuje dokładny zbiór dostępnych
warstw. Parser tekstu otrzymuje ten zbiór przy każdym zapytaniu i nie może
wybrać parametru, dla którego nie opublikowano wyników.

\section{Bezpieczna sekwencja aktywacji}

\begin{enumerate}[nosep]
  \item definicja i aliasy;
  \item pobieranie bieżące;
  \item historia i RangeAwareBackfillBridge;
  \item audyt pokrycia, jednostek i duplikatów;
  \item opcjonalna rola cechy;
  \item trening quick i brama jakości;
  \item trening full;
  \item powierzchnia przestrzenna i publikacja.
\end{enumerate}

\end{document}
